\documentclass[11pt]{article}
% ======================================================================
%  examstyle.tex — EPFL-style exam layout for the weekly Time Series exams
%  Shared by every Exam_Week_NN(.tex / _Solutions.tex). Compiles with pdflatex.
% ======================================================================
\usepackage[utf8]{inputenc}
\usepackage[T1]{fontenc}
\usepackage[english]{babel}
\usepackage[a4paper,margin=2.4cm]{geometry}
\usepackage{amsmath,amssymb,amsthm,mathtools}
\usepackage{bm}
\usepackage{enumitem}
\usepackage{array}

\setlength{\parindent}{0pt}
\setlength{\parskip}{0.5em}
\emergencystretch=3em
\allowdisplaybreaks[1]

% ---------- math macros (same names as the rest of the project) ----------
\newcommand{\E}{\mathbb{E}}
\DeclareMathOperator{\Var}{Var}
\DeclareMathOperator{\Cov}{Cov}
\DeclareMathOperator{\Corr}{Corr}
\DeclareMathOperator*{\argmin}{arg\,min}
\DeclareMathOperator{\MSE}{MSE}
\DeclareMathOperator{\trace}{tr}
\newcommand{\Reals}{\mathbb{R}}
\newcommand{\Ints}{\mathbb{Z}}
\newcommand{\Comp}{\mathbb{C}}
\newcommand{\Nat}{\mathbb{N}}
\newcommand{\Prob}{\mathbb{P}}
\newcommand{\Tr}{^{\mathsf{T}}}
\newcommand{\conj}[1]{\overline{#1}}
\newcommand{\abs}[1]{\left\lvert #1 \right\rvert}
\newcommand{\norm}[1]{\left\lVert #1 \right\rVert}
\newcommand{\ind}{\mathbf{1}}
\newcommand{\eps}{\varepsilon}
\newcommand{\diff}{\nabla}
\newcommand{\acvs}{\gamma}
\newcommand{\acf}{\rho}
\newcommand{\pacf}{\alpha}
\newcommand{\sdf}{\mathcal{S}}
\newcommand{\Bsh}{B}
\newcommand{\iid}{\overset{\text{iid}}{\sim}}

\setlist[enumerate]{leftmargin=2em,topsep=2pt,itemsep=4pt}
\setlist[itemize]{leftmargin=1.6em,topsep=2pt,itemsep=2pt}

\newcounter{exo}

% ---------- exam cover header :  \examheader{exam title}{duration} ----------
\newcommand{\examheader}[2]{%
  \thispagestyle{empty}
  \begin{center}
    {\Large\bfseries Time Series}\\[3pt]
    Spring Semester 2026, EPFL\\
    \'Ecole Polytechnique F\'ed\'erale de Lausanne\\
    Prof.\ Dr.\ Sofia C.\ Olhede\\[7pt]
    \hrule height 0.8pt
    \vspace{9pt}
    {\large\bfseries Practice Exam --- #1}\\[4pt]
    {\normalsize Duration: #2 \quad$\cdot$\quad No calculator \quad$\cdot$\quad Answer on paper}\\
    \vspace{8pt}
    \hrule height 0.8pt
  \end{center}
  \vspace{14pt}
  \begin{tabular}{@{}p{8.2cm}@{}p{8cm}@{}}
    Name:\enspace\hrulefill & Surname:\enspace\hrulefill
  \end{tabular}\\[14pt]
  SCIPER (Student Card Number):\enspace\hrulefill
  \vspace{16pt}
}

% ---------- points table (fixed: 4 problems) ----------
\newcommand{\pointstable}{%
  \begin{center}
  \renewcommand{\arraystretch}{1.5}
  \begin{tabular}{|p{5cm}|p{4cm}|}
    \hline
    \textbf{Exercise} & \textbf{Points}\\\hline
    1 & \\\hline
    2 & \\\hline
    3 & \\\hline
    4 & \\\hline
    \textbf{Total} & \\\hline
  \end{tabular}
  \end{center}
}

% ---------- problem heading (exam: one problem per page) ----------
\newcommand{\problem}[1]{%
  \clearpage
  \stepcounter{exo}%
  \begin{center}\textbf{\large Problem \theexo\ (#1 points)}\end{center}
  \vspace{4pt}
}
% ---------- answer space: fill the statement page + one blank answer page ----------
% Put \answerpage at the END of each problem so every exercise gets its own page(s).
\newcommand{\answerpage}{\par\vfill\newpage\null\newpage}

% ---------- solutions document ----------
\newcommand{\solheader}[1]{%
  \begin{center}
    {\Large\bfseries Time Series --- Solutions}\\[3pt]
    {\large Practice Exam --- #1}\\[2pt]
    EPFL, MATH-342 \quad$\cdot$\quad Prof.\ S.\ C.\ Olhede\\[5pt]
    \hrule height 0.8pt
  \end{center}
  \vspace{10pt}
}
\newcommand{\solproblem}[1]{%
  \bigskip
  \stepcounter{exo}%
  {\large\bfseries Problem \theexo\ (#1 points)}\par\nobreak\vspace{2pt}
}
% Rappel de l'énoncé avant la solution (dans les corrigés)
\newcommand{\statementstart}{\par\vspace{1pt}{\bfseries\itshape Statement.}\par\nobreak\begingroup\small\itshape}
\newcommand{\statementend}{\par\endgroup\vspace{5pt}\hrule\vspace{6pt}{\bfseries Solution.}\par\nobreak\vspace{2pt}}

% --- local macros not in examstyle ---
\newcommand{\Herm}{^{\mathsf{H}}}
\DeclareMathOperator{\diag}{diag}
\begin{document}

\examheader{Week 9: Multivariate Time Series \& VAR}{60 minutes}
\pointstable
\begin{center}\itshape Justify every answer. Throughout, $\eps_t$ is a mean-zero white noise of variance $\sigma^2$.\end{center}

% ---------------------------------------------------------------
\problem{5}
We work with bivariate series $X_t=(X_t^{(1)},X_t^{(2)})\Tr$. Recall the matrix
autocovariance $\Gamma_\tau=\Cov(X_{t+\tau},X_t)$ (\emph{the lag $\tau$ sits in
the first argument}), with $(j,k)$ entry the cross-covariance
$\acvs_\tau^{(j,k)}=\Cov(X_{t+\tau}^{(j)},X_t^{(k)})$, and the
cross-correlation $\acf_\tau^{(j,k)}=\acvs_\tau^{(j,k)}/\sqrt{\acvs_0^{(j,j)}\acvs_0^{(k,k)}}$.

\textbf{(a)} State the symmetry relation linking $\Gamma_{-\tau}$ to $\Gamma_\tau$, prove it from the definition, and explain in one sentence why the off-diagonal cross-covariance $\acvs_\tau^{(1,2)}$ need \emph{not} be an even function of $\tau$ (unlike a scalar ACVS). \hfill(1.5 pts)

\textbf{(b)} Let $\{U_t^{(1)}\}$ and $\{U_t^{(2)}\}$ be two jointly stationary, mutually independent zero-mean series with $\Var(U_t^{(j)})=\sigma_j^2$, and let $\{\phi_t\}$ be a scalar mean-zero white noise of variance $\sigma_\phi^2$, independent of both. Define the common-shock model
\[
Y_t^{(1)}=U_t^{(1)}+\phi_t,\qquad Y_t^{(2)}=U_t^{(2)}+\phi_t .
\]
Compute the cross-covariance $\acvs_\tau^{(1,2)}$ and the cross-correlation $\acf_\tau^{(1,2)}$ between $\{Y^{(1)}\}$ and $\{Y^{(2)}\}$ at \emph{all} lags $\tau$, and identify at which lag(s) the shared shock injects correlation. \hfill(2 pts)

\textbf{(c)} Specialise to $\sigma_1^2=3$, $\sigma_2^2=1$, $\sigma_\phi^2=1$ and report the numerical value of $\acf_0^{(1,2)}$ and of $\acf_\tau^{(1,2)}$ for $\tau\neq0$. \hfill(1.5 pts)
\answerpage

% ---------------------------------------------------------------
\problem{5}
\textbf{(The delay process and its cross-spectrum.)} Fix an integer delay $\nu$ and a gain $\alpha>0$. Let $\{X_t^{(2)}\}$ be a zero-mean stationary series with autocovariance $\acvs_\tau^{(2,2)}$, spectral density $\sdf_2(f)=\sum_\tau\acvs_\tau^{(2,2)}e^{-2\pi if\tau}$ and variance $\sigma_2^2=\acvs_0^{(2,2)}$. Define a second channel as a scaled, delayed copy corrupted by noise,
\[
X_t^{(1)}=\alpha\,X_{t-\nu}^{(2)}+\eps_t ,\qquad t\in\Ints,
\]
where $\{\eps_t\}$ (variance $\sigma^2$) is white noise uncorrelated with $\{X_t^{(2)}\}$ at all lags.

\textbf{(a)} Show that the cross-covariance is a \emph{shift} of the autocovariance of $X^{(2)}$, namely $\acvs_\tau^{(1,2)}=\alpha\,\acvs_{\tau-\nu}^{(2,2)}$, and deduce the cross-correlation $\acf_\tau^{(1,2)}$ at all lags in terms of $\acf^{(2,2)}$. \hfill(1.5 pts)

\textbf{(b)} Compute the cross-spectrum $S_{1,2}(f)=\sum_\tau\acvs_\tau^{(1,2)}e^{-2\pi if\tau}$ in closed form, and read off its \textbf{amplitude} $\abs{S_{1,2}(f)}$ and its \textbf{phase} $\arg S_{1,2}(f)$. Why is the phase \emph{linear} in $f$, and what is its slope? \hfill(2 pts)

\textbf{(c)} Define the coherence $r_{1,2}(f)=S_{1,2}(f)/\sqrt{S_{1,1}(f)\,S_{2,2}(f)}$. Using $S_{1,1}(f)=\alpha^2\sdf_2(f)+\sigma^2$ and $S_{2,2}(f)=\sdf_2(f)$, give $\abs{r_{1,2}(f)}$, and show that in the \emph{noiseless} limit $\sigma^2\to0$ the magnitude-squared coherence equals $1$ at every frequency. \hfill(1.5 pts)
\answerpage

% ---------------------------------------------------------------
\problem{5}
\textbf{(VAR stability and the VAR(1) second-order structure.)}

\textbf{(a)} Consider the bivariate $\mathrm{VAR}(2)$ process $X_t=\Phi_1X_{t-1}+\Phi_2X_{t-2}+\eps_t$ with
\[
\Phi_1=\begin{pmatrix}0.5 & 0\\[2pt] 0.3 & 0.7\end{pmatrix},\qquad
\Phi_2=\begin{pmatrix}0 & 0\\[2pt] 0.1 & -0.1\end{pmatrix}.
\]
Form the reverse characteristic polynomial $\det\bigl(I_2-\Phi_1 z-\Phi_2 z^2\bigr)$, exploit the lower-triangular structure to \emph{factor} it, find all its roots, and decide whether the VAR(2) is stable. State the companion-matrix ($F$) form of the stability test and verify it gives the same conclusion. \hfill(2.5 pts)

\textbf{(b)} Now take the \emph{diagonal} $\mathrm{VAR}(1)$ process $X_t=\Phi\,X_{t-1}+\eps_t$ with
\[
\Phi=\diag\!\bigl(\tfrac12,\tfrac13\bigr),\qquad
\{\eps_t\}\sim\mathrm{WN}(0,I_2).
\]
Write its $\mathrm{MA}(\infty)$ representation, and compute the stationary covariance $\Gamma_0$ and the lag-$1$ matrix $\Gamma_1$ in closed form (exact fractions). \hfill(1.5 pts)

\textbf{(c)} For the process of part (b), what is the coherence $r_{1,2}(f)$ between the two channels at every frequency, and why? \hfill(1 pt)
\answerpage

% ---------------------------------------------------------------
\problem{5}
\textbf{(Revision of Week 8: the tapered periodogram --- bias kernel \& unbiasedness.)} Data $X_1,\dots,X_n$ ($T=\{1,\dots,n\}$, $\Delta=1$) are observed from a zero-mean stationary process with spectral density $\sdf(f)$. A \emph{data taper} is a sequence $\{h_t\}_{t\in T}$ (extended by $h_t=0$ off $T$); the tapered periodogram is
\[
\hat\sdf^{(p)}_h(f)=\abs{J_h(f)}^2 ,\qquad
J_h(f)=\sum_{t\in T}h_t\,X_t\,e^{-2\pi itf},
\]
and $H(f)=\sum_{t}h_t e^{-2\pi itf}$ denotes the discrete Fourier transform of the taper.

\textbf{(a)} Show that the expectation of the tapered periodogram is the frequency convolution
\[
\E\bigl[\hat\sdf^{(p)}_h(f)\bigr]=\int_{-1/2}^{1/2}\sdf(f')\,\abs{H(f-f')}^2\,df' .
\]
(\emph{Hint:} expand $\abs{J_h}^2$ as a double sum, take expectations to get $\sum_\tau w_\tau\acvs_\tau e^{-2\pi if\tau}$ with $w_\tau=\sum_t h_t h_{t+\tau}$, then identify $w$ as an autocorrelation of the taper and apply the convolution theorem.) \hfill(2 pts)

\textbf{(b)} Suppose $\{X_t\}$ is white noise of variance $\sigma^2$ (so $\sdf\equiv\sigma^2$). Using part (a) and Parseval's relation $\int_{-1/2}^{1/2}\abs{H(f')}^2\,df'=\sum_t h_t^2$, prove that $\hat\sdf^{(p)}_h(f)$ is \emph{unbiased} for $\sdf(f)$ at every $f$ and every $n$ \textbf{iff} $\norm{h}_2^2=\sum_t h_t^2=1$. \hfill(1.5 pts)

\textbf{(c)} Take $n=3$ and the un-normalised triangular weights $g=(g_1,g_2,g_3)=(1,2,1)$. Find the constant $c$ so that $h_t=c\,g_t$ is an admissible taper (i.e.\ $\norm h_2^2=1$), and contrast, in one or two sentences, what tapering versus \emph{multitapering} each fix about the raw periodogram. \hfill(1.5 pts)
\answerpage

\end{document}
